\begin{abstract}

Model-based control and Reinforcement Learning address different limitations in
robotic locomotion. Model Predictive Control (MPC) uses an explicit dynamics
model to optimize future motion and handle physical constraints, but its
performance can degrade when the model does not capture the full behavior of the
system. Reinforcement Learning (RL) can learn behaviors directly through
interaction, making it suitable for compensating for effects that are difficult
to model. However, policies trained from scratch often require extensive reward
design and provide little insight into the resulting control strategy. This
thesis investigates a hybrid approach in which a learned policy augments a
nominal MPC controller instead of replacing it. The model-based component
provides the main control action, while the policy learns corrections aimed at
improving closed-loop performance. The framework is implemented in JAX using
MPX to parallelize the MPC computations required during training and is
evaluated in the MuJoCo physics engine. Its performance is assessed against both
model-based and learning-based baselines, with particular attention to tracking
accuracy, robustness, and generalization across different commands and
operating conditions.

\end{abstract}